\chapter{移居者权益与义务}\label{ux79fbux5c45ux8005ux6743ux76caux4e0eux4e49ux52a1}

本世界以``资源充足、公共供给、各取所需''为基本经济前提；社会运行不以工资与固定岗位为核心，而以项目制工坊、公共服务网络与自治协作维持长期稳定。移居者制度的目标只有两个：\textbf{把人安全、平等地接入公共供给}；\textbf{把公共系统的安全边界讲清楚并可执行}。我们拒绝用``少数人审批权''去管理移民，也拒绝用``全系统自动化''去假装没有风险------因此本章采用\textbf{分布式身份、可审计但不可滥用、权利默认生效}的原则。

\section{身份认定与社会保障}\label{ux8eabux4efdux8ba4ux5b9aux4e0eux793eux4f1aux4fddux969c}

\subsection{身份层级（从低到高）}\label{ux8eabux4efdux5c42ux7ea7ux4eceux4f4eux5230ux9ad8}

\begin{itemize}
\item
  \textbf{访客身份（Visitor）}
\item
  适用：短期停留、考察、探亲。
\item
  权利：公共医疗急救、基础住宿与餐食、公共交通、公共空间使用。
\item
  限制：不得接入高权限科研设施、不得持有高功率制造权限、不得调用公共机器人编队。
\item
  \textbf{临时常住（Provisional Resident）}
\item
  适用：移居初期过渡。
\item
  默认时长：\textbf{180天}（可因学习/康复/司法程序延长）。
\item
  权利：完整公共医疗、子女教育、基本住房配置、个人AI与生活机器人服务配额、法律援助、语言与适应课程。
\item
  义务：完成安全与公民课程（见13.2）、遵守机器人/算力/能源边界。
\item
  \textbf{常住（Resident）}
\item
  适用：已稳定生活、可参与中长期项目。
\item
  获得条件：完成180天过渡 + 无重大安全违规 + 完成一次社会贡献结项（劳动/科研/艺术/志愿任一）。
\item
  权利：与公民等同的公共供给（除少数安全敏感权限），可发起公共项目。
\item
  \textbf{公民（Citizen）}
\item
  获得条件（基准）：\textbf{连续居住2年} + 完成至少2个公共项目结项 + 通过``公共安全边界''复核。
\item
  权利：参与制度迭代投票/审议、参与跨世界协作代表团抽选池、可申请安全敏感设施的受控权限。
\end{itemize}

\begin{quote}
注：期限与条件是``默认基准''，其目的不是设门槛，而是把\textbf{安全训练、协作习惯、责任链}做实。

\end{quote}

\subsection{身份认定方式（不依赖少数人权力）}\label{ux8eabux4efdux8ba4ux5b9aux65b9ux5f0fux4e0dux4f9dux8d56ux5c11ux6570ux4ebaux6743ux529b}

\begin{itemize}
\item
  \textbf{分布式身份凭证（DID式）}：移居者获得一份加密身份根证书；由多个独立节点共同签发与更新，避免单点机构``卡脖子''。
\item
  \textbf{最小化披露}：日常只验证``你有资格做这件事''，不暴露你的全部身份信息。
\item
  \textbf{可审计但去标识}：公共资源调用（住房、算力、机器人、医疗等）写入不可篡改日志，但默认去标识；只有触发安全事件才进入受控揭示流程。
\item
  \textbf{线下应急身份}：灾害/断网时使用离线令牌保证医疗与避难权利不断供。
\end{itemize}

\subsection{社会保障包（默认即生效）}\label{ux793eux4f1aux4fddux969cux5305ux9ed8ux8ba4ux5373ux751fux6548}

移居者从``临时常住''起默认获得以下保障，不与``贡献多少''绑定：

\begin{itemize}
\item
  \textbf{住房}：基础居住单元按``\textbf{约120 m³/人}''配置（示例：45㎡ × 2.7m层高 ≈ 121.5 m³；本手册空间基准）。允许单亲家庭/照护需求上浮。
\item
  \textbf{医疗}：免费公共医疗 + 机器人医疗体系；含心理健康与创伤支持。
\item
  \textbf{教育}：学校制 + 项目制工坊，无分轨；儿童与成人都可进入再学习路径。
\item
  \textbf{生活协助}：配备生活服务型机器人/人形服务者，强制要求``可硬件断开、不可诱导''的交互边界（本项目既定安全设定）。
\item
  \textbf{基础算力与AI}：每人有基础个人AI额度，用于学习、创作、生活管理；高能耗训练/大规模仿真需走公共项目申报（见13.2）。
\item
  \textbf{能源权利}：本世界采用既定的人均能源预算模型（项目内设定：\textbf{16.8万kWh/人/年，总能耗结构保持}；来源：项目《能源消耗结构》讨论）。移居者从常住起与本地居民同等适用。
\end{itemize}

\begin{center}\rule{0.5\linewidth}{0.5pt}\end{center}

\section{社会贡献形式（劳动、科研、艺术、社会志愿服务）}\label{ux793eux4f1aux8d21ux732eux5f62ux5f0fux52b3ux52a8ux79d1ux7814ux827aux672fux793eux4f1aux5fd7ux613fux670dux52a1}

我们不以``岗位---工资''组织社会，但仍需要\textbf{可预测的公共维护与创新产出}。因此采用``贡献即参与治理''的框架：贡献不是交换生存权，而是进入更大权限与更深协作网络的通行证。

\subsection{最低贡献要求（基准）}\label{ux6700ux4f4eux8d21ux732eux8981ux6c42ux57faux51c6}

\begin{itemize}
\item
  \textbf{临时常住}：完成公共安全与适应课程 + 至少一次社区协作（合计建议\textbf{40小时/180天}，可拆分）。
\item
  \textbf{常住}：每年建议完成\textbf{120小时/人/年}的社会贡献（四类任选，可混合）。
\item
  \textbf{豁免}：未成年人、重病/残障、主要照护者可豁免或改为轻量贡献（例如知识分享、陪伴支持）。
\end{itemize}

\begin{quote}
这不是``劳动换口粮''，而是用低门槛保证：系统不会被大量``完全不参与的人群''拖入不可治理状态。

\end{quote}

\subsection{四类贡献的可计量口径}\label{ux56dbux7c7bux8d21ux732eux7684ux53efux8ba1ux91cfux53e3ux5f84}

\begin{itemize}
\item
  \textbf{劳动类}（维护、建造、运营）
\item
  例：公共空间维护、应急演练、能源与供水设施巡检、机器人街区``断电地板''维护。
\item
  计量：工时 + 风险等级 + 影响范围（由项目看板公开）。
\item
  \textbf{科研类}（从实验到工程化）
\item
  例：能源系统优化、机器人安全机制、卫星与航天工程。
\item
  计量：可复现成果（代码/图纸/实验记录）+ 同行评审 + 落地贡献。
\item
  能源口径：对高能耗项目实行``能耗预算''透明制。项目内已有跨世界航天协作能耗规则基准（例如联合航天研发：各组\textbf{300 kWh/人/年}；特定调研组\textbf{3000 kWh/人/年}；来源：项目《国际通讯/航天协议》讨论），科研项目可参照此类``明确到人均''的分摊方式。
\item
  \textbf{艺术与文化类}（公共精神生活基础设施）
\item
  例：公共展演、叙事档案、社区节律设计、夜生活综合体的节目与策展。
\item
  计量：公共可访问性 + 参与人数 + 持续性（一次性爆发不等于长期供给）。
\item
  \textbf{社会志愿服务}（照护、调解、教育支持）
\item
  例：新移民适应陪伴、语言伙伴、儿童项目辅导、冲突调解支持。
\item
  计量：服务时长 + 受益反馈 + 风险/保密等级。
\end{itemize}

\subsection{贡献与权限的对应（``能力---责任''对齐）}\label{ux8d21ux732eux4e0eux6743ux9650ux7684ux5bf9ux5e94ux80fdux529bux8d23ux4efbux5bf9ux9f50}

\begin{itemize}
\item
  基础生活权利不挂钩；但以下权限与贡献挂钩：

  \begin{itemize}
  \item
    高能耗算力调用（训练大型模型、全城仿真）
  \item
    高级制造权限（高能量密度材料、关键基础设施部件）
  \item
    机器人编队与公共空间大规模调度权
  \item
    安全敏感区域/数据访问
  \end{itemize}
\item
  对应机制：公开项目申请 + 预算（能耗/空间/材料）+ 多节点签发权限 + 全程日志审计。

\end{itemize}

\begin{center}\rule{0.5\linewidth}{0.5pt}\end{center}

\section{科研成果归属}\label{ux79d1ux7814ux6210ux679cux5f52ux5c5e}

在公有经济下，``归属''不应等同于``占有''，而应等同于\textbf{可追溯的署名权、可执行的开放规则、可控的安全边界}。我们采用三层规则：

\subsection{默认：公共知识许可}\label{ux9ed8ux8ba4ux516cux5171ux77e5ux8bc6ux8bb8ux53ef}

\begin{itemize}
\item
  \textbf{默认开放}：公共资源（公共算力、公共实验室、公共材料）产出的科研成果默认进入``公共知识库''。
\item
  \textbf{署名与贡献链}：所有贡献者按可追溯记录署名；署名带来的是声誉、话语权与项目优先级，而不是私人垄断。
\item
  \textbf{可复现要求}：关键结论必须附带可复现实验/仿真条件或替代验证路径。
\end{itemize}

\subsection{例外：安全与伦理受控开放}\label{ux4f8bux5916ux5b89ux5168ux4e0eux4f26ux7406ux53d7ux63a7ux5f00ux653e}

以下成果进入``受控开放库''，只向合格人员开放，并强制风险评估：

\begin{itemize}
\item
  可显著提升攻击能力的机器人控制、破坏关键基础设施的方法
\item
  高危生物/化学/材料合成路径
\item
  可用于强监控、强操控的行为诱导系统
  受控并不意味着永久封存：每一项受控都有\textbf{定期解密评审周期}，以免``安全''为名造成知识垄断。
\end{itemize}

\subsection{私人创作与商业交换（有限存在）}\label{ux79c1ux4ebaux521bux4f5cux4e0eux5546ux4e1aux4ea4ux6362ux6709ux9650ux5b58ux5728}

\begin{itemize}
\item
  纯个人资源、纯私人时间完成的作品允许保留更强的个人处置权（尤其艺术作品）。
\item
  但只要接入公共关键资源（公共算力/实验室/关键材料），就必须至少回馈一个层级的公共版本（例如开源核心、封装应用可差异化）。
\item
  目的：鼓励个人创造，同时防止公共系统被``借壳私有化''。
\end{itemize}

\begin{center}\rule{0.5\linewidth}{0.5pt}\end{center}

\section{国际矛盾下的移民处理办法}\label{ux56fdux9645ux77dbux76feux4e0bux7684ux79fbux6c11ux5904ux7406ux529eux6cd5}

我们假设``跨世界/跨集团矛盾''长期存在，因此移民政策必须在冲突中仍然可执行。原则只有四条：\textbf{生存权不断供、个体责任不连坐、程序正义可复核、冲突隔离可落地}。

\subsection{身份与保护状态}\label{ux8eabux4efdux4e0eux4fddux62a4ux72b6ux6001}

\begin{itemize}
\item
  冲突期新增三种保护标记（不改变你原本身份层级）：
\item
  \textbf{中立保护（Neutral Protection）}：来自冲突方但未参与敌对行动者，享受完整生活供给与人身保护。
\item
  \textbf{人道庇护（Humanitarian Asylum）}：存在迫害/战乱风险者，优先安置与心理/医疗支持。
\item
  \textbf{受限接入（Restricted Access）}：存在高概率安全风险者，限制其进入关键设施与高权限算力/机器人调度，但基础生活权利不受影响。
\end{itemize}

\subsection{冲突期的处理流程（可审计、可申诉）}\label{ux51b2ux7a81ux671fux7684ux5904ux7406ux6d41ux7a0bux53efux5ba1ux8ba1ux53efux7533ux8bc9}

\begin{itemize}
\item
  \textbf{快速风险分流（48小时内）}：只做最低必要判断：是否需要隔离关键权限、是否需要医疗与庇护。
\item
  \textbf{受控调查（30天内）}：多节点调查组给出结论与理由，所有证据链进入受控审计。
\item
  \textbf{申诉与复核}：移居者可申请一次独立复核；复核组与初审组必须节点隔离。
\item
  \textbf{最小限制原则}：除非有明确证据，限制只能落在``权限''而非``生存''。
\end{itemize}

\subsection{迁徙与遣返}\label{ux8fc1ux5f99ux4e0eux9063ux8fd4}

\begin{itemize}
\item
  \textbf{不强制遣返到迫害风险地}（人道底线）。
\item
  若移居者自愿返回，系统需提供最低限度的旅程保障与资产/署名权保全。
\item
  若发生严重安全犯罪：

  \begin{itemize}
  \item
    先司法审理，再决定隔离、改造性项目参与或跨域移交；
  \item
    禁止``群体性驱逐''。
  \end{itemize}
\end{itemize}

\subsection{社区层面的冲突隔离设计}\label{ux793eux533aux5c42ux9762ux7684ux51b2ux7a81ux9694ux79bbux8bbeux8ba1}

\begin{itemize}
\item
  在城市空间上提供``低摩擦隔离带''：中立居住区、公共调解站、跨文化工坊，以减少群体对立被放大。
\item
  在机器人与算力层面启用``冲突期策略''：关键基础设施的机器人只接受更窄的指令集；高风险模型训练自动降级到更严格的能耗与权限预算（与本项目既定的``硬件断开、可审计、不可诱导''安全边界一致）。
\end{itemize}
